% The listings are read from the repository at compile time, not copied:
%   04_Study/4.4_App-Sampling/scripts/googleplay-scraper/v007/  sample.py, config.json, requirements.txt
\chapter{Sampling Script}
\label{app:script}
The script in Listing~\ref{lst:sample} drew the sample of
Section~\ref{sec:results-sample}. It needs Python~3 and the open-source
parser \texttt{google-play-scraper} (version 1.2.7), and it runs in three
steps. First, it searches Google Play for every query of
Table~\ref{tab:keywords} and keeps up to the first 30 apps of each result page.
Second, it reads the store page of each app it found and of the preinstalled
app, and stores each page in a cache file, so that an interrupted run can
resume. Third, it ranks the apps by their German review count and walks down
the ranking. It takes every app that the screening file marks as eligible and
stops when the sample is full. When it reaches an app that has not been
screened yet, it stops and names that app, so the screening is done by hand in
rank order. The search results and the store pages are saved, and a second run
reuses them, so the screening does not change the candidates.

\begin{table}[H]
\centering
\begin{singlespace}
\footnotesize
\caption[Search queries]{Search queries. An asterisk marks a word that is the
same in English and German and was searched once. The last row holds the
queries added by the check with well-known apps.}
\label{tab:keywords}
\begin{tabularx}{\textwidth}{@{}l>{\raggedright\arraybackslash}X>{\raggedright\arraybackslash}Xl@{}}
\toprule
\textbf{Category} & \textbf{English} & \textbf{German} & \textbf{Source} \\
\midrule
% Generated by 04_Study/4.4_App-Sampling/scripts/googleplay-scraper/v007/tables.py; do not edit by hand.
Cycling & cycling & Radfahren & Hutton et al. \\
Diet & diet, eating & Diät, Ernährung & Hutton et al. \\
Exercise & exercise, workouts* & Training & Hutton et al. \\
Heart rate & heartrate, heart rate & Herzfrequenz, Puls & Hutton et al. \\
Mood & mood, happiness & Stimmung, Glück & Hutton et al. \\
Running & walking, running & Gehen, Laufen & Hutton et al. \\
Sleep & sleep & Schlaf & Hutton et al. \\
Step count & steps & Schritte & Hutton et al. \\
Weight & weight, body fat & Gewicht, Körperfett & Hutton et al. \\
Reproductive health & female health, women's health, reproductive health, fertility, ovulation, menstrual, menstruation*, period, contraception, birth control, pregnancy, pregnant & Frauengesundheit, reproduktive Gesundheit, Fruchtbarkeit, Eisprung, Periode, Verhütung, Schwangerschaft, schwanger & Malki et al. \\
Devices and scales & fitness tracker, smartband*, scale & Fitnesstracker, Waage & Known-app check \\
\bottomrule
% the generated rows end with \bottomrule
\end{tabularx}
\end{singlespace}
\end{table}

The script reads its settings from a JSON file (Listing~\ref{lst:config}):
the queries, the storefront (\texttt{region}), the interface language, the
number of results kept per query (\texttt{depth}), the sample size, the
package name of the preinstalled app, the screening file and the output
folder. The paths are relative to the root of the project repository, from
which the script is run:
\texttt{python 04\_Study/4.4\_App-Sampling/scripts/googleplay-scraper/v007/sample.py}
\texttt{04\_Study/4.4\_App-Sampling/scripts/googleplay-scraper/v007/config.json}.
It writes \texttt{search.json}, \texttt{pages.jsonl}, \texttt{candidates.csv}
and \texttt{sample.csv} to the output folder.

\lstinputlisting[style=jsoncode, caption={Settings of the run of 21 November 2026; the queries after the second are left out.}, label={lst:config}, linerange={1-4,57-67}]{../04_Study/4.4_App-Sampling/scripts/googleplay-scraper/v007/config.json}

\lstinputlisting[style=pythoncode, caption={Sampling script (\texttt{sample.py}).}, label={lst:sample}]{../04_Study/4.4_App-Sampling/scripts/googleplay-scraper/v007/sample.py}
